\section{The numerical argument on a good minimal model}
\label{sec:geometry}

We prove the fiber-space conclusion assuming that the total space has
a good klt minimal model. The upper bound comes from the geometric
generic fiber. For the lower bound, we construct a mixed intersection
that detects the ample class from the base and kills both terms in
the canonical comparison.

\begin{theorem}[The good-model case]\label{thm:good-model-case}
Let $f\colon X\to Y$ be a surjective morphism with connected fibers
between smooth connected projective complex varieties. Let $F$ be its
smooth geometric generic fiber and suppose $\kappa(F)=0$. Assume
that $X$ has a projective good klt minimal model. If $H$ is an ample
Cartier divisor on $Y$ and $m_0$ is a positive integer such that
\begin{equation}\label{eq:geometric-numerical-hypothesis}
 m_0K_X-f^*H\in\PE(X),
\end{equation}
then $\kappa(X)=\dim Y$.
\end{theorem}

We use the good-model convention and effective comparison
\eqref{eq:intro-comparison} stated in the introduction. This is the
$K_X$-negative convention of \cite[Definitions~3.6.1 and~3.6.7]{BCHM2010}.

\subsection{The upper bound from the generic fiber}

We first account for the use of $\kappa(F)=0$. It bounds every
pluricanonical image of $X$, before any model is chosen.

\begin{lemma}[The geometric generic fiber bound]\label{lem:easy-addition}
Let $f\colon X\to Y$ be a surjective morphism with connected fibers
between smooth connected projective complex varieties. If its smooth
geometric generic fiber $F$ satisfies $\kappa(F)=0$, then
$\kappa(X)\leq\dim Y$.
\end{lemma}

\begin{proof}
Put $K=\C(Y)$ and $\eta=\operatorname{Spec}K$. Generic smoothness
and Stein factorization give a smooth geometrically integral generic
fiber $X_\eta$ \cite[III, Sections~10--11]{Hartshorne1977}.
The cotangent sequence along this fiber gives
\begin{equation}\label{eq:canonical-generic-restriction}
 \omega_X|_{X_\eta}
 \simeq\omega_{X_\eta/K}\otimes_K\ell,
 \qquad \ell=\omega_Y|_\eta.
\end{equation}
Here $\ell$ is a one-dimensional $K$-vector space. The factor from
the base therefore cancels in ratios of pluricanonical sections.

Fix $m>0$ such that $H^0(X,mK_X)\ne0$, and choose a basis
$s_0,\ldots,s_N$ with $s_0\ne0$. A nonzero section is nonzero at
the generic point of $X$, so it stays nonzero on $X_\eta$ and after
extension to the geometric generic fiber. By
\eqref{eq:canonical-generic-restriction}, the restricted system on
$F$ is a nonzero subsystem of $|mK_F|$. Its image has dimension zero,
since $\kappa(F)=0$.

Let $\phi_m$ be the rational map of the complete system $|mK_X|$,
and let $T$ be the closure of the image of
$(f,\phi_m)\colon X\dashrightarrow Y\times\mathbb P^N$.
Its function field is
\[
 L=K(s_1/s_0,\ldots,s_N/s_0)\subseteq\C(X).
\]
The dimension of its generic image over $Y$ is
$\operatorname{trdeg}_K L$. This dimension is unchanged by algebraic
extension of $K$. Geometric integrality of $X_\eta$ identifies that
extended image with the image of the same section ratios on $F$.
Thus $\operatorname{trdeg}_K L=0$, and $\dim T=\dim Y$.
Projection to $\mathbb P^N$ gives $\dim\phi_m(X)\leq\dim Y$.
Taking the maximum over all nonempty pluricanonical systems proves
the lemma. If every system is empty, the inequality is immediate.
\end{proof}

\subsection{A curve class that detects the base}

The lower bound requires a different ingredient. The following lemma
constructs a numerical test from a globally generated divisor on $V$.
All its factors come from $V$, which is why exceptional errors vanish.
Its strict positivity uses the independent morphism to $Y$.

\begin{lemma}[A mixed intersection test]\label{lem:mixed-positivity}
Let $q\colon W\to V$ be a birational morphism from a smooth
projective complex $d$-fold to a normal projective variety.
Let $D$ be a globally generated Cartier divisor on $V$, and let
$g\colon V\to Z\subseteq\mathbb P^N$ be the morphism of $|D|$
onto its image. Write $k=\dim Z$, and let $A$ be a very ample
Cartier divisor on $V$. Suppose $h\colon W\to Y$ is a surjective
morphism to a smooth projective variety with $y=\dim Y>k$.
For every ample Cartier divisor $H$ on $Y$, the numerical one-cycle
class
\[
 C=(q^*D)^k(q^*A)^{d-k-1}
\]
satisfies the following properties:
\begin{enumerate}
 \item $B\cdot C\geq0$ for every pseudo-effective real divisor
 class $B$ on $W$;
 \item $E\cdot C=0$ for every $q$-exceptional rational divisor $E$;
 \item $q^*D\cdot C=0$ and $h^*H\cdot C>0$.
\end{enumerate}
\end{lemma}

\begin{proof}
The inequality $k<y\leq d$ makes all exponents nonnegative; a
product with no factors has its usual meaning. We prove first that
the test is nonnegative, then that it has the two required vanishing
properties, and finally that it detects the ample divisor from $Y$.

\smallskip\noindent
\emph{Nonnegativity.}
The divisors $q^*D$ and $q^*A$ are globally generated. If $B$ is
effective, choose their members successively so that no member
contains an irreducible component of the preceding intersection
on $B$. Global generation permits all the finitely many required
avoidances. The resulting Cartier intersections form an effective
zero-cycle, possibly empty. Hence $B\cdot C\geq0$
\cite[Chapter~2]{Fulton1998}. Extend by linearity to effective real
divisors and by continuity on $N^1(W)_\R$ to its closed effective
cone. This proves (1).

\smallskip\noindent
\emph{Vanishing on exceptional and semiample classes.}
For a $q$-exceptional divisor the Cartier projection formula gives
\begin{equation}\label{eq:exceptional-annihilation}
 E\cdot(q^*D)^k(q^*A)^{d-k-1}
   =q_*E\cdot D^kA^{d-k-1}=0;
\end{equation}
see \cite[Proposition~2.3(c)]{Fulton1998}.
Indeed, every component of $E$ maps to codimension at least two,
so $q_*E=0$ as a divisor cycle. This reasoning is valid on the
possibly singular variety $V$.
Also $D=g^*\cO_Z(1)$, and $k+1$ general hyperplanes miss the
$k$-dimensional image $Z$. Consequently $(q^*D)^{k+1}=0$.
For $k=0$, the morphism $g$ is constant and $\cO_V(D)$ is trivial,
which gives the same conclusion. These observations prove (2) and
the first assertion of (3).

\smallskip\noindent
\emph{Strict positivity.}
Choose $b>0$ such that $bH$ is very ample. Use this divisor to
embed $Y$, and use $A$ to embed $V$.
There is a point $w\in W$ at which $q$ is an isomorphism onto a
smooth open subset of $V$ and
\[
 \operatorname{rank}d(gq)_w=k,
 \qquad\operatorname{rank}dh_w=y,
 \qquad\operatorname{rank}dq_w=d.
\]
Each condition holds on a dense open subset. For the rank statements,
the generic differential rank in characteristic zero equals the
dimension of the image, by separability of the function-field
extension. We choose $w$ in the intersection of these open subsets.

The differentials of pulled-back hyperplanes through $gq(w)$ span
a $k$-dimensional subspace $U\subseteq T_w^*W$. Choose $k$ of them
whose differentials form a basis of $U$. The hyperplanes through
$h(w)$ supply a $y$-dimensional subspace. Since $y>k$, one such
hyperplane pulls back to a divisor with differential outside $U$.
Finally, hyperplanes from the embedding by $A$ supply all cotangent
directions at $w$, because $q$ is a local isomorphism there. Choose
$d-k-1$ of them completing the chosen differentials to a basis of
$T_w^*W$. These $d$ divisors meet transversely at $w$.

To use this local intersection in the global intersection number,
perturb the hyperplanes slightly in their complex parameter spaces.
The transverse point persists for every sufficiently small
perturbation, by the implicit function theorem. Tuples whose
pullbacks meet properly at every successive step form a nonempty
Zariski open subset of the product of the hyperplane parameter
spaces. Nonemptiness follows by successively avoiding the finitely
many components already present, using global generation; openness
follows from upper semicontinuity of fiber dimension for the
projective universal intersections. This parameter space is
irreducible, so the open subset is dense also in the complex
analytic topology. It therefore meets the perturbation neighborhood.

The resulting proper global intersection is an effective zero-cycle
that contains a transverse point of multiplicity one. It follows that
\[
 h^*(bH)\cdot(q^*D)^k(q^*A)^{d-k-1}>0.
\]
Division by $b$ proves (3). Only one cotangent direction from $h$
outside $U$ was needed; the two morphisms $gq$ and $h$ need not
factor through each other.
\end{proof}

\subsection{Sections and the numerical contradiction}

We now combine the two lemmas. The good model supplies actual
sections on $X$. The mixed test forces their image to have at least
the dimension of $Y$.

\begin{proof}[Proof of Theorem~\ref{thm:good-model-case}]
If $Y$ is a point, then $F=X$ and $\kappa(F)=0$ is the required
conclusion. Hence assume $d=\dim X\geq y=\dim Y>0$.

Let $V$ be a good klt minimal model of $X$. Resolve the closure of
the graph of $X\dashrightarrow V$ projectively
\cite{Hironaka1964}, obtaining a smooth projective common resolution
$X\xleftarrow{p}W\xrightarrow{q}V$. By the good-model comparison,
compatible canonical divisors satisfy
\begin{equation}\label{eq:canonical-comparison}
 p^*K_X=q^*K_V+E,\qquad E\geq0,\qquad q_*E=0.
\end{equation}
Pulling back to a further common resolution preserves the effectivity
and target-exceptionality of this error.

Choose $r>0$ such that $D=rK_V$ is Cartier and globally generated,
and let $g\colon V\to Z\subseteq\mathbb P^N$ be its morphism onto
its image. Put $k=\dim Z$. Since $K_X$ is Cartier,
\eqref{eq:canonical-comparison} shows that
$rE=p^*(rK_X)-q^*D$ is an integral Cartier divisor on $W$.
The two maps whose dimensions we will compare appear in
Figure~\ref{fig:common-resolution}.

\begin{figure}[ht]
\centering
\begin{tikzcd}[column sep=large,row sep=large]
 & W \arrow[dl,"p"'] \arrow[dr,"q"] & \\
 X \arrow[d,"f"'] && V \arrow[d,"g"] \\
 Y && Z
\end{tikzcd}
\caption{The common resolution carries the two morphisms $h=fp$ to
$Y$ and $gq$ to $Z$.
The divisor factors defining the mixed class are pulled back from
$V$, so the class annihilates the
$q$-exceptional error. Its pairing with $h^*H$ detects the ample
class from the separate base $Y$.}
\label{fig:common-resolution}
\end{figure}

Pull back sections of $D$ by $q$ and multiply by the canonical
section of the effective Cartier divisor $rE$. We obtain an injection
\begin{equation}\label{eq:section-transfer}
 H^0(V,D)\hookrightarrow H^0(W,rp^*K_X)=H^0(X,rK_X).
\end{equation}
The equality follows from $p_*\cO_W=\cO_X$ and the sheaf projection
formula for the proper birational map to the normal variety $X$
\cite[III, Corollary~11.4 and its proof; II, Exercise~5.1(d)]{Hartshorne1977}.
Multiplication by the section of $rE$ leaves section ratios unchanged
on the complement of its support. Thus the subsystem in
\eqref{eq:section-transfer} has image dimension $k$, and
\begin{equation}\label{eq:canonical-image-lower-bound}
 \kappa(X)\geq k.
\end{equation}

Suppose, for a contradiction, that $k<y$. Choose a very ample
Cartier divisor $A$ on $V$ and apply
Lemma~\ref{lem:mixed-positivity} with
\[
 h=fp,\qquad C=(q^*D)^k(q^*A)^{d-k-1}.
\]
The canonical comparison and the vanishing assertions of the lemma
give
\begin{equation}\label{eq:canonical-test-zero}
 p^*K_X\cdot C
 =\frac1r q^*D\cdot C+E\cdot C=0.
\end{equation}
Pullback by $p$ preserves pseudo-effectivity: on smooth varieties
it takes effective real divisors to effective real divisors and
induces a continuous linear map on numerical divisor spaces.
Therefore the pullback of
\eqref{eq:geometric-numerical-hypothesis}, paired with $C$, yields
\[
 0\leq(m_0p^*K_X-h^*H)\cdot C=-h^*H\cdot C<0,
\]
a contradiction. Hence $k\geq y$. Equation
\eqref{eq:canonical-image-lower-bound} and
Lemma~\ref{lem:easy-addition} now give
$y\leq k\leq\kappa(X)\leq y$, proving the theorem.
\end{proof}

The proof includes $k=0$, $d-k-1=0$ and relative dimension zero.
Effectivity of $E$ supplies the section injection, while
target-exceptionality makes its mixed pairing vanish. These are
separate uses of the canonical comparison, and both are needed.
